\clearpage\thispagestyle{empty}
\pdfbookmark[0]{Quellen zu Gauß’ Meridianbeobachtungen des Neptun (August 1847)}{neptune1847-sources}
\begin{center}\textbf{Quellen zu Gauß’ Meridianbeobachtungen des Neptun (August 1847)}\end{center}
\noindent Diese KI-Übertragung ins heutige Deutsch folgt den Beobachtungen auf S.~472 der \href{https://gdz.sub.uni-goettingen.de/id/PPN236007467}{\textit{Werke}, Band~VI} (Göttingen, 1874). Carl Friedrich Gauß machte die drei Beobachtungen am 12., 13. und 14. August 1847 in Göttingen und datierte den Bericht auf den 15. August. Der Erstdruck steht in den von Heinrich Christian Schumacher herausgegebenen \textit{Astronomischen Nachrichten}, Band~XXVI, Nr.~610, \href{https://api.digitale-sammlungen.de/iiif/image/v2/bsb10049451_00085/full/1800,/0/default.jpg}{Sp.~153--154}. Die \href{https://api.digitale-sammlungen.de/iiif/image/v2/bsb10049451_00088/full/1800,/0/default.jpg}{Schlusszeile auf Sp.~160} nennt den 30. August 1847. Der Katalog der Bayerischen Staatsbibliothek datiert den gebundenen Band auf 1848.
\par\smallskip\noindent Der Erstdruck enthält Gauß’ Unterschrift, kürzt die Überschriften ab und verwendet andere Zeitzeichen und Dezimaltrennzeichen. Tabelle und Verfasserangabe folgen hier der Werkausgabe. KI-Übertragung und KI-Satz.
\clearpage
